% Edited conversion of source p. 20.
\ReportHeading
  {Нелінійне деформування функціонально-градієнтних матеріалів, складові яких мають ефект пам’яті форми}
  {П.~О. Стеблянко}
  {Інститут механіки ім. С.~П. Тимошенка НАН України}
  {}

Узагальнено визначальні нелінійні співвідношення термопластичності для функціонально-градієнтних матеріалів (ФГМ), складові яких мають ефект пам’яті форми.

Для коефіцієнтів визначальних співвідношень, що описують залежність характеристик ФГМ від виду навантаження під час деформування за межами пружності, проведено верифікацію параметрів матеріалу на основі відомих експериментальних даних для функціональних матеріалів. Тензор деформацій у визначальних рівняннях подано як суму тензорів пружної та пластичної деформацій.

Фізичні рівняння в задачах термопластичності записано у вигляді
\begin{equation}
  \sigma_{ij}=a_{ijkl}\varepsilon_{kl}+b_{ij},
\end{equation}
де $a_{ijkl}=\mathrm{const}$ у пружному випадку або залежить від параметрів попереднього термопластичного деформування. Для ФГМ і функціональних матеріалів ці параметри, своєю чергою, можуть залежати від матеріальних координат тіла та часу.

Для металокерамічних ФГМ фізико-механічні властивості матеріалу є функціями об’ємних часток його складових і можуть бути подані за правилом сумішей
\begin{equation}
  P=P_mV_m+P_cV_c,
  \qquad V_m+V_c=1,
\end{equation}
де $V_m$ і $V_c$ --- об’ємні частки металевої та керамічної складових відповідно.

Розглянуто лінійну та кубічну інтерполяції зміни складу ФГМ уздовж однієї з координат. Узагальнено визначальні співвідношення, що описують залежність характеристик матеріалу від виду навантаження під час деформування за межами пружності, та верифіковано параметри моделі за відомими експериментальними даними для функціональних матеріалів.

Розроблено метод розв’язування стаціонарних задач деформування тіл обертання з матеріалів із пам’яттю форми. Для підвищення точності в ньому застосовано сплайн-функції.
